\section{Datenbank, Ausführungsmaschine und Anbindung an Sprachmodelle}\label{sec:maschine}

\subsection{Die Datenbank}\label{sec:datenbank}

Für einen Bestand aus vielen Dokumenten liest
\begin{lstlisting}
python -m qam.tools.readout_db index
\end{lstlisting}
alle exportierten Akten in eine lokale SQLite-Datenbank. Dabei werden Eintrags- und Aktenhashes geprüft und die Prämissenkanten in beiden
Richtungen angelegt. Neuindizieren erneuert diese Kanten, lässt aber gespeicherte Belege, den Rechencache und ausdrücklich registrierte
Ersetzungen unverändert. \texttt{find} sucht Einträge mit Frame, Status und Prämissen; \texttt{premises HASH} liefert alles, worauf
ein Eintrag transitiv ruht; \texttt{impact HASH} alle späteren Einträge und Belege, die ihn verwendet haben, auch als fehlende Prämisse;
\texttt{register-replacement ALT NEU --reason} verknüpft eine ausdrückliche Korrektur, ohne den alten Eintrag oder alte Belege zu ändern.
Wer eine Aussage wiederverwenden will, ruft sie so mit Status, Frame und Prämissen ab; wer sie korrigiert, erfährt, welche Schlüsse
betroffen sind.

\subsection{Aufträge und Belege: MML und MMA}

Rechnungen werden als \emph{Aufträge} in einer kleinen Auftragssprache (MML) gestellt und von einem Ausführer (MMA) bearbeitet. Ein Aufruf
besteht aus zwei JSON-Dateien: einem \emph{Kontext}, der die beteiligten QAM-Objekte und Guards deklariert, und einem \emph{Auftrag}, der eine
Anweisung mit Argumenten nennt:
\begin{lstlisting}
python -m qam.tools.readout_db mma KONTEXT.json AUFTRAG.json --summary
\end{lstlisting}
MMA prüft zuerst die Guards von Kontext, Auftrag und Objekten und die verknüpften Prämissen, dann führt es aus und speichert einen
\emph{Beleg}: den vollständigen Auftrag, den Ausgang (\texttt{continue}, \texttt{guard}, \texttt{reject}), das Ergebnis oder die Gründe der
Sperre, einen Kontexthash, Fingerabdrücke der beteiligten Implementierung und gegebenenfalls die Prämissenspur. Der Beleg ist durch einen
SHA-512-Hash identifiziert. Ein Zähler in der Datenbank zählt verschiedene Belege je Ausgang; derselbe Beleg zählt nur einmal.

Ein gespeicherter Beleg lässt sich zweifach prüfen. \texttt{receipt HASH} prüft die gespeicherten Hashes; \texttt{receipt HASH --replay}
rechnet den Auftrag mit dem gespeicherten Kontext und der aktuellen Implementierung ohne Cache neu und meldet getrennt, ob Beleg und
Ergebnis übereinstimmen. Ein geänderter Implementierungshash kann den Beleg ändern, auch wenn das Ergebnis gleich bleibt. Der Rechencache
schließt Argumente, Kontext, Prämissenhashes und Implementierungsversionen in seinen Schlüssel ein; Guards werden geprüft, bevor der Cache
befragt wird.

Die Anweisungen sind bewusst eng: Es gibt keine freie Programmausführung, sondern eine feste Liste. Beispiele:
\begin{itemize}[nosep]
\item \texttt{READOUT} führt ein Rezept aus einer zugelassenen Liste aus, mit geprüften Argumentbereichen.
\item \texttt{SEARCH\_ROUTE} durchsucht einen festen Katalog von Operationen bis zu einer Schrittgrenze und speichert neben den Kandidaten
die ausgeschlossenen Wege mit Grund und die offenen Guards jedes Kandidaten. Findet sie keinen Weg ohne offene Annahme, endet sie mit
\texttt{guard}.
\item \texttt{RUN\_TAPE} führt eine deklarierte Bandmaschine aus (Abschnitt~\ref{sec:band}).
\end{itemize}

\paragraph{Stapel.} Bis zu 64 unabhängige Aufträge mit eindeutigen Namen laufen in einem Stapel auf ein bis vier Prozessen. Abhängigkeiten
(\texttt{depends\_on}, \texttt{requires\_match}) werden vor dem Start auf Zyklenfreiheit geprüft; gesperrte Nachfolger erscheinen als
\texttt{blocked} mit \texttt{blocked\_by} (Abschnitt~\ref{sec:ebeneB}). Der Stapel erzeugt einen gehashten Bericht in Eingabereihenfolge,
der mit den Einzelbelegen in einer Transaktion gespeichert wird. Berichte können neu gerechnet (\texttt{replay-batch}), verglichen
(\texttt{compare-batches}) und ohne Neurechnung auf vollständige Verweise geprüft werden (\texttt{audit-batch}).

\subsection{Plan und Ausführung am Beispiel der Bandmaschine}\label{sec:band}

Die Anweisung \texttt{RUN\_TAPE} zeigt die Rollenteilung der Einleitung in ihrer einfachsten Form. Der \emph{Plan} besteht aus zwei
QAM-Objekten des Kontexts: einem \texttt{operator} mit Zuständen, Haltezustand und einer deterministischen binären Übergangstabelle, und
einem \texttt{state} mit Band, Kopfposition und Steuerzustand. Jeder Tabelleneintrag legt zu Steuerzustand und gelesenem Bit das zu
schreibende Bit, die Kopfbewegung \(-1,0,+1\) und den Folgezustand fest. Der Auftrag nennt beide Objekte und setzt die Betriebsart
(\texttt{fixed} oder \texttt{grow}, bei der rechts eine leere Zelle angefügt wird), Grenzen für Schritte und Zellen (höchstens je 1024) und
optional eine Ablaufspur.

Das Ergebnis unterscheidet zwei Fragen, die leicht verwechselt werden:
\begin{itemize}[nosep]
\item \emph{Durfte der Auftrag ausgeführt werden?} Das ist der MMA-Ausgang, etwa \texttt{continue}.
\item \emph{Wie endete die Rechnung?} Das ist der Stoppgrund der Bandmaschine: \texttt{halted} (regulärer Halt), \texttt{stuck} (keine
Regel für Zustand und Bit), \texttt{boundary} (Schritt über eine feste Bandgrenze, der ganze Schritt wird verworfen), \texttt{step\_limit}
oder \texttt{cell\_limit} (Budget erschöpft).
\end{itemize}
Ein erschöpftes Budget ist ausdrücklich kein Nichthalten. Die Maschine behauptet nicht, die Rechnung würde nie enden; sie meldet, dass sie
innerhalb des Budgets nicht endete, und gibt Band, Kopf und Zustand zurück, mit denen ein neuer Lauf mit größerem Budget fortsetzen kann.
Das ist die Turing-Grenze in ihrer praktischen Form.

Derselbe Plan wird von mehreren Ausführenden bearbeitet: vom Python-Adapter, von einem Wolfram-Language-Skript, von einer Qiskit-Fassung,
die jeden Schritt als reversible Schaltung über alle Entscheidungsfälle baut, und von zwei SystemVerilog-Steuerungen. Eine
Vergleichssuite führt beide Steuerungen auf 209 Referenzfällen gegen MMA aus und vergleicht Stoppgründe, Taktzahlen und eine
Beispielspur; optional werden Logikzellen synthetisiert. Die binäre Übergangstabelle und das Bandwachstum sind Modellentscheidungen; geprüft
wird die ausführbare Übergangssemantik, nicht eine Deutung der Maschine.

\subsection{Warum ein Band: eine Maschine im mathematischen Träger}\label{sec:traeger}

Die Bandmaschine ist kein beliebig gewähltes Testbeispiel. Sie ist das ausführbare Gegenstück zu einer Lesart, die der Forschungstext
des Projekts~\cite{companion} verfolgt: ob die mathematische Struktur selbst eine Rechenmaschine trägt, also eine virtuelle Maschine,
die nicht auf Hardware, sondern in einem mathematischen Träger läuft. Der Text gibt dafür eine lokale Übergangssemantik auf Zuständen
\(\Sigma=(\tau,h,q,\rho,\chi)\) aus Bandausschnitt, Kopfposition, Steuerzustand, Auslesung und Abschlusszertifikat an. Die Primitive
\texttt{FORWARD} und \texttt{HOPF} bewegen den Kopf, \texttt{PROJECT}, \texttt{COMPARE}, \texttt{DIFF} und \texttt{EXPORT\_BOOL} liefern
Verzweigungsprädikate, \texttt{COLLAPSE}, \texttt{CLOSE} und \texttt{HCSEL} ändern den Steuerzustand, \texttt{REJECT} hält an.

Die Rolle des Bands erhält die \glqq optimale Linie\grqq{} des zulässigen spinoriellen Transports. Das ist die \emph{spinorielle
Bandhypothese} (STH): Diese Linie ist ein fortsetzbarer Lese- und Schreibträger, entlang dessen endliche Bandausschnitte verlängert
werden können, ohne die Zulässigkeit zu verletzen. Der Grund für diese Wahl ist strukturell. Transport entlang einer Linie ist sequentiell
und eindimensional, und das Band ist in Turings Modell der einfachste Speicher mit lokalem Zugriff; ein Transportschritt ist dann eine
Kopfbewegung. Das Band ist zugleich als stabilster Bestandteil vorgesehen, um den herum Steuerung, Verzweigung und Rückkehr aufgebaut
werden.

Die zweite Hypothese trägt Zuses Namen. Die Z3 kannte keine bedingte Verzweigung; Rojas hat gezeigt, dass sie trotzdem im Prinzip
universell ist, indem Verzweigung und indirekte Adressierung arithmetisch nachgebildet werden, so universell wie ein Rechner mit
beschränktem Adressraum und um den Preis sehr langer Programme~\cite{rojas98}. Die \emph{bedingte Z3-Hypothese} (Z3H) überträgt das:
Sobald Verzweigung, beständiger Speicher und Rückkehr auf dem spinoriellen Träger verfügbar sind, kann der Stapel MML--MMA--MCP einen
gewöhnlichen sequentiellen Steuerkern nachbilden. Die \glqq Tricks\grqq{}, mit denen eine solche Maschine universell würde, sind also benannt.

Bewiesen ist die Universalität nicht, und der Forschungstext sagt das ausdrücklich. Er isoliert einen bedingten Weg mit fünf
Verpflichtungen in fester Reihenfolge: die Steuerschicht setzt sich über benachbarte Fenster fort (\(U_Q\)), die Verzweigung ist lokal
eindeutig ausführbar (\(U_B\)), die Rückkehr über den HC-gewählten Abschluss ist fortsetzbar (\(U_L\)), das Band ist global fortsetzbar
(\(U_T\)), und ein Referenzkern wird explizit simuliert (\(U_S\)). In QAM ist \glqq der Träger ist Turing-vollständig\grqq{} damit eine Aussage
mit Status \st{conditional} und fünf offenen Bedingungen, und keine Rechnung der Maschine darf eine davon stillschweigend streichen.

Die Ausführung mit \texttt{RUN\_TAPE} prüft etwas anderes, nämlich die Übergangssemantik des Plans. Eine deterministische Maschine mit
binärem Alphabet und einem einseitig unbeschränkten Band ist im idealisierten Sinn universell; mit höchstens 1024 Schritten und Zellen je
Auftrag ist der Ausführer aber, wie die Z3 bei Rojas, nur im beschränkten Sinn universell. Die Identifikation des Bands mit dem spinoriellen
Träger bleibt eine Voraussetzung, die kein Lauf erfüllt.

Hier schließt sich der Bogen zu Turing. Wäre der Träger universell, dann wäre für seine Programme das Halteproblem unentscheidbar. Dass
\texttt{step\_limit} und \texttt{cell\_limit} als eigene Stoppgründe geführt und nie als \glqq hält nicht\grqq{} gelesen werden, ist deshalb
keine Bequemlichkeit, sondern folgt aus der Grenze, die das Modell selbst erreichen soll.

\subsection{Anbindung an Sprachmodelle über MCP}\label{sec:mcp}

Damit ein Sprachmodell mit der Akte arbeiten kann, ohne sie verändern zu können, stellt QAM seine Abfragen über das Model Context
Protocol~\cite{mcp} bereit. Es gibt drei Server mit abgestuften Rechten:
\begin{center}
\small
\begin{tabular}{@{}p{0.2\textwidth}p{0.42\textwidth}p{0.3\textwidth}@{}}
\toprule
Server & Werkzeuge & Rechte\\
\midrule
schreibgeschützt & Einträge suchen, Belege auflisten und abrufen (mit Spur oder Zusammenfassung), Kontextauswahl, Stapelberichte auflisten, lesen, vergleichen und prüfen & Container ohne Netz, Datenbank nur lesend eingebunden\\
vollständig & zusätzlich Index neu bauen, Einzelaufträge und gespeicherte Stapel ausführen & schreibt Belege in die Datenbank\\
Stapel & gespeicherte oder direkt übergebene Stapel ausführen, Berichte neu rechnen, HDL-Vergleich & schreibt Belege\\
\bottomrule
\end{tabular}
\end{center}
Der schreibgeschützte Server läuft in einem Docker-Container mit \texttt{--read-only}, \texttt{--network none} und nur lesend
eingebundener Datenbank. Alle Server führen nur feste Module mit geprüften Argumenten und einer Zeitgrenze aus; Auftragsdateien müssen in zugelassenen
Verzeichnissen liegen. Keiner der Server kann einen Status anheben, eine Akte umschreiben oder einen Beleg löschen.

Die Zusammenfassung eines Belegs nennt geprüften Hash, Anweisung, Ausgang, Fehler, Verpflichtungen und offene Guards. Ein Modell, das
eine Frage zu einem gespeicherten Ergebnis beantwortet, ruft den Beleg ab, prüft ihn per Hash und nennt die offenen Guards, statt sich auf
ein früheres Gespräch zu verlassen. Das Sprachmodell ist wieder das Rechenwerk: Es formuliert Aufträge und liest Belege, aber was als
Ergebnis gilt, entscheidet die Maschine nach dem Plan.
